\documentclass[11pt]{article}
\usepackage[margin=1in]{geometry}
\usepackage[T1]{fontenc}
\usepackage{lmodern,amsmath,amssymb,amsthm}
\usepackage[hidelinks]{hyperref}
\usepackage{microtype}
\newtheorem{theorem}{Theorem}
\setlength{\parskip}{0.35em}
\setlength{\parindent}{0pt}
\title{Tightness of contact Godbillon--Vey forms:\\
 a classical argument for Calegari's conditional question}
\author{Alec Kriebel\\\small Independent researcher\\
\small ORCID: \href{https://orcid.org/0009-0001-9320-500X}{0009-0001-9320-500X}}
\date{5 October 2026}
\begin{document}
\maketitle
\begin{abstract}
If a global defining form $\alpha$ of a cooriented taut $C^2$ foliation
on a closed three-manifold satisfies $d\alpha=\alpha\wedge\omega$
with a smooth contact form $\omega$, then $\ker\omega$ is tight.
The proof combines a special case of Dathe--Khoule's 2012 general
affine-deformation criterion with the Eliashberg--Thurston tightness
neighborhood theorem and smooth Gray stability. A finite-interval
smoothing argument allows $\alpha$ to be only $C^1$.
This is a classical corollary, not a claim of a new contact construction
or a first historical solution. Priority remains unresolved.
\end{abstract}

\section{The conditional question and the precise claim}
The target is Question 13.2 on printed p.~29 of Calegari's actual
2002 author preprint, dated 8 September, Version 0.78 \cite{C}.
Its antecedents are inherited from Question 13.1: $\mathcal F$ is a
minimal taut $C^2$ foliation of an atoroidal three-manifold $M$,
and $\mathrm{gv}(\mathcal F)[M]\ne0$. One additionally supposes
that a choice of global defining form $\alpha$ makes the associated
one-form $\omega$ contact, with the convention
\begin{equation}
T\mathcal F=\ker\alpha,\qquad d\alpha=\alpha\wedge\omega.
\label{connection}
\end{equation}
The requested conclusion is that the contact structure $\ker\omega$
is tight. We use the closed, oriented interpretation supplied by a
nonzero evaluation on the ordinary fundamental class. This is an
interpretation of that evaluation, not a quoted source-wide convention;
no boundary or noncompact version is asserted.

Question 13.1 separately asks whether the defining form can be chosen
so that $\omega\wedge d\omega$ has one weak sign everywhere.
Neither that existence question nor existence of a contact $\omega$
is established here. The final AMS chapter published in 2003 has DOI
\href{https://doi.org/10.1090/pspum/071/2024640}{10.1090/pspum/071/2024640};
its full text has not been read, and equality with the 2002 wording
is not presumed.

\begin{theorem}\label{main}
Let $M$ be a closed oriented smooth three-manifold, and let
$\mathcal F$ be a cooriented taut $C^2$ foliation without spherical
leaves. Suppose $\alpha$ is a global nowhere-zero $C^1$ one-form
defining $\mathcal F$, and $\omega$ is a smooth contact one-form
satisfying \eqref{connection}. Then $\ker\omega$ is tight in the
usual smooth contact-topological sense.
\end{theorem}

Minimality in the source excludes spherical leaves. Atoroidality,
minimality, and the numerical Godbillon--Vey value play no further
role in this implication. The global defining form supplies the
coorientation. The conclusion therefore applies to the specified
2002 conditional question under the smooth contact-form interpretation.

\section{Proof by a finite contact path}
We use two classical inputs. First, every smooth contact structure
sufficiently $C^0$ close to the tangent planes of a taut $C^2$
foliation without spherical leaves on a closed three-manifold is
tight. This Eliashberg--Thurston neighborhood statement is explicitly
recorded in Vogel \cite[p.~42]{V11} and Dathe--Rukimbira
\cite[Proposition 3.3, p.~5]{DR}; the latter refers to p.~50 of
\emph{Confoliations} \cite{ET}. We rely on those actual primary-paper
statements, rather than claiming a full reading of the book.
For negative contact structures apply the same result after reversing
the orientation of the relevant component. Second, smooth Gray
stability identifies every member of a smooth family of smooth contact
structures on a closed manifold by an isotopy
\cite[Theorem 2.9, p.~2451]{V16}. Tightness is consequently constant
on such a path.

\begin{proof}
Differentiating \eqref{connection} gives
\begin{equation}
0=d^2\alpha=d\alpha\wedge\omega-\alpha\wedge d\omega
=-\alpha\wedge d\omega.
\label{mixed}
\end{equation}
This calculation is valid when $\alpha$ is only $C^1$: $d^2\alpha=0$
distributionally, $\alpha\wedge\omega$ is $C^1$, and its product
rule yields the continuous form $\alpha\wedge d\omega$.
A continuous form that vanishes distributionally vanishes pointwise.
For each constant real $s$, define $\beta_s=\omega+s\alpha$. Then
\begin{align}
\beta_s\wedge d\beta_s
&=\omega\wedge d\omega
 +s(\omega\wedge d\alpha+\alpha\wedge d\omega)
 +s^2\alpha\wedge d\alpha \notag\\
&=\omega\wedge d\omega.\label{volume}
\end{align}
Thus every $\beta_s$ is contact, with the same sign. Moreover,
\begin{equation}
\ker\beta_s=\ker(\alpha+s^{-1}\omega)\longrightarrow
\ker\alpha\quad\text{uniformly in }C^0\quad(s\longrightarrow+\infty).
\label{limit}
\end{equation}
Compactness and the positive minimum of $|\alpha|$ justify this
uniform plane-field convergence. Choose a finite $S>0$ for which
$\ker\beta_S$ lies in the tightness neighborhood just described.

If $\alpha$ is smooth, \eqref{volume} gives a smooth contact path on
$[0,S]$. Its endpoint is tight, so smooth Gray stability proves the
claim at $s=0$.

For $C^1$ $\alpha$, keep this same finite $S$ and keep $\omega$
unchanged. Approximate $\alpha$ in $C^1$ by a smooth one-form $a$,
using smooth charts and a partition of unity, and consider
\begin{equation}
\lambda_s=\omega+s a,\qquad 0\le s\le S.
\label{smoothpath}
\end{equation}
Here is a uniform contact bound. Choose a metric and, component by
component, an auxiliary orientation making
$\omega\wedge d\omega=c(x)dV$ with $c\ge c_0>0$.
On $M\times[0,S]$, bound both $|\beta_s|$ and $|d\beta_s|$ by $B$.
If the one-form and exterior-derivative errors of \eqref{smoothpath}
are at most $\delta$, their contact-volume error is at most
$2B\delta+\delta^2$. Choose $\delta$ with this bound below $c_0/2$;
the approximation can make the errors at most $\delta$ uniformly
because $S$ is fixed. Every $\lambda_s$ is then contact.
Also $\ker\lambda_S$ remains inside the $C^0$ tightness neighborhood
by taking the approximation sufficiently small. Smooth Gray stability
along \eqref{smoothpath} transports its tightness to
$\ker\lambda_0=\ker\omega$.
\end{proof}

The approximation $a$ need not define a foliation or satisfy
\eqref{connection}. Only closeness on a fixed finite interval is used.
Gray stability is never applied to a $C^1$ path, to $s=\infty$, or
to a foliated endpoint. Reversing an auxiliary orientation does not
change tightness of the underlying contact planes.

\section{Classical attribution and the limits of the note}
Dathe--Khoule's general affine criterion
\cite[Definition 2.1 and Theorem 2.2, pp.~102--103]{DK} already
covers the contact mechanism. In dimension three, for an integrable
defining form $\alpha$ and a positive contact form $\omega$, its
criterion is
\begin{equation}
Q=\alpha\wedge d\omega+\omega\wedge d\alpha\ge0.
\label{DKcriterion}
\end{equation}
Their definition permits $C(t)=1$, $B(t)=t$, giving
$\alpha+t\omega$. In the present case each term of $Q$ is zero,
and their expansion specializes to
\[
(\alpha+t\omega)\wedge d(\alpha+t\omega)
=t^2\omega\wedge d\omega\qquad(t>0).
\]
After positive scaling, these are precisely the planes of
$\omega+s\alpha$ with $s=1/t$. The criterion assumes integrability,
not closedness of $\alpha$, and allows zero mixed term; no strictly
positive first variation is needed. Negative sign is handled by
orientation reversal.

The original 2012 statement was read in an author-uploaded article
transcription; its PDF pixels were not obtained. It does not itself
print a $C^1$ threshold. The later coauthor preprint
\cite[\S6, Theorem 6.2, pp.~32--33]{KMNW} explicitly attributes the
$C^1$ criterion to the 2012 work. Its normalized-equality display
after equation (62) omits the general nonnegative mixed term and
should be an inequality. Equality is correct for $Q=0$, the only
case used here; \eqref{volume} independently verifies it.
No unrelated main result of that 2025 preprint is endorsed.

Theorem~\ref{main} is our stated corollary of prior affine,
Eliashberg--Thurston, and Gray results. The near-foliation/isotopy
tightness argument already occurs with a closed defining form in
Dathe--Rukimbira \cite[Proposition 3.4]{DR}; that narrower statement
alone does not cover nonclosed $\alpha$.
No literal statement of the exact tightness answer to Calegari's
Question 13.2 was located in the complete transcribed 2012 article
or the relevant 2025 section. This bounded absence is not firstness
evidence. The final 2003 chapter, the full \emph{Confoliations} book,
the final 2011 Dathe--Rukimbira article, and other identified literature
remain unread. Historical priority, novelty of any explicit application
or regularity refinement, and present-day open status are unresolved.

The portable support includes the longer current candidate, which
separately treats $C^1$ $\omega$ using an explicitly defined embedded
$C^2$-disk criterion and tightness of sufficiently $C^1$-close smooth
contact forms. Those qualified extras are not claims of this short
theorem, and are not identified with every low-regularity definition
of tightness. Exact symbolic and rational diagnostics check wedge
identities, sign conventions, smoothing corrections, and finite-interval
margins. They neither prove the imported global contact-topology
theorems nor certify priority.

AI tools were used extensively in analysis, drafting, source comparison,
and verification. No conventional human peer review or refereeing has
occurred. Independent AI audits are not human peer review.

\clearpage
\begin{thebibliography}{9}\small
\bibitem{C} D. Calegari, \emph{Problems in foliations and laminations
of 3-manifolds}, Version 0.78 (8 September 2002), Definition 1.1,
p.~1; Questions 13.1--13.2, p.~29.
\href{https://arxiv.org/abs/math/0209081v1}{arXiv:math/0209081v1}.

\bibitem{DK} H. Dathe and C. Khoule, \emph{Sur les d\'eformations
d'un feuilletage de codimension 1 en structures de contact},
African Diaspora J. Math. \textbf{13}(2) (2012), 100--107.
Definition 2.1 and Theorem 2.2, pp.~102--103.
\href{https://www.researchgate.net/publication/258228980_Sur_les_deformations_d\%27un_feuilletage_de_codimension_1_en_structures_de_contact}{Author-uploaded original article transcription}.

\bibitem{ET} Y. Eliashberg and W. P. Thurston, \emph{Confoliations},
University Lecture Series \textbf{13}, AMS (1998), p.~50 cited through
the explicit statements in \cite{V11,DR}; full book not read.

\bibitem{V11} T. Vogel, \emph{Rigidity versus flexibility for tight
confoliations}, Geom. Topol. \textbf{15} (2011), 41--121;
tightness neighborhood, p.~42.
\href{https://doi.org/10.2140/gt.2011.15.41}{doi:10.2140/gt.2011.15.41}.

\bibitem{V16} T. Vogel, \emph{On the uniqueness of the contact structure
approximating a foliation}, Geom. Topol. \textbf{20} (2016), 2439--2573;
Theorem 2.9, p.~2451.
\href{https://doi.org/10.2140/gt.2016.20.2439}{doi:10.2140/gt.2016.20.2439}.

\bibitem{DR} H. Dathe and P. Rukimbira, \emph{Contact deformations of
closed 1-forms on torus bundles over the circle}, author preprint (2008),
Propositions 3.3--3.4, p.~5.
\href{https://arxiv.org/abs/0812.3389v1}{arXiv:0812.3389v1}.

\bibitem{KMNW} C. Khoule, M. Manso, A. Ndiaye and K. War,
\emph{$C^0$-Contact Anosov flows}, arXiv:2503.00454v1 (1 March 2025),
\S6, Theorem 6.2, pp.~32--33; only this relevant criterion and
attribution are used.
\href{https://arxiv.org/abs/2503.00454v1}{arXiv:2503.00454v1}.
\end{thebibliography}
\end{document}
